\documentclass[12pt, oneside]{article}

% ---------- Packages ----------
\usepackage[T1]{fontenc}
\usepackage{mathpazo}
\usepackage[margin=1.25in]{geometry}
\usepackage{amsmath, amssymb, amsthm}
\usepackage{microtype}
\usepackage{setspace}
\usepackage{titlesec}
\usepackage{hyperref}
\usepackage{xcolor}
\usepackage{csquotes}
\usepackage{enumitem}
\usepackage{fancyhdr}
\usepackage{abstract}

% ---------- Hyperref ----------
\hypersetup{
  colorlinks=true,
  linkcolor=black,
  citecolor=black,
  urlcolor=black,
  pdftitle={Constraint Closure and Its Failures: Projection, Entropy, and the Degeneration of Form in AI Systems and Written Discourse},
  pdfauthor={Flyxion},
}

% ---------- Typography ----------
\setstretch{1.15}
\setlength{\parskip}{0.45em}
\setlength{\parindent}{1.5em}

\titleformat{\section}{\large\scshape}{\thesection.}{0.75em}{}
\titleformat{\subsection}{\normalsize\itshape}{\thesubsection.}{0.5em}{}

% ---------- Header / Footer ----------
\pagestyle{fancy}
\fancyhf{}
\fancyhead[L]{\small\scshape Constraint Closure and Its Failures}
\fancyhead[R]{\small\thepage}
\renewcommand{\headrulewidth}{0.3pt}

% ---------- Theorem environments ----------
\theoremstyle{plain}
\newtheorem{theorem}{Theorem}[section]
\newtheorem{proposition}[theorem]{Proposition}
\newtheorem{lemma}[theorem]{Lemma}
\newtheorem{corollary}[theorem]{Corollary}

\theoremstyle{definition}
\newtheorem{definition}[theorem]{Definition}
\newtheorem{example}[theorem]{Example}

\theoremstyle{remark}
\newtheorem{remark}[theorem]{Remark}

% ---------- Abstract style ----------
\renewcommand{\abstractnamefont}{\scshape\small}
\renewcommand{\abstracttextfont}{\normalfont\small}
\setlength{\absleftindent}{1.5em}
\setlength{\absrightindent}{1.5em}

% ---------- Math macros ----------
\newcommand{\R}{\mathbb{R}}
\newcommand{\N}{\mathbb{N}}
\newcommand{\proj}{\Pi}
\newcommand{\energy}{\mathcal{E}}
\newcommand{\config}{\mathcal{X}}
\newcommand{\obs}{\mathcal{O}}
\newcommand{\traj}{\gamma}
\newcommand{\attractor}{\mathcal{A}}
\newcommand{\prior}{\mathcal{P}}

% ==========================================================
\begin{document}
% ==========================================================

\begin{titlepage}
\vspace*{2cm}
\begin{center}
{\LARGE\scshape Constraint Closure and Its Failures}\\[1em]
{\large Projection, Entropy, and the Degeneration of Form\\
in AI Systems and Written Discourse}\\[3em]
{\normalsize Flyxion}\\[0.5em]
{\small Independent Researcher}\\[2em]
{\small\today}
\end{center}

\vspace{2em}
\begin{abstract}
\noindent
This essay develops a unified formal framework for analyzing two apparently distinct phenomena:
the accidental release of proprietary AI system code and the degeneration of argumentative essay
form under generative pressure. Both are shown to be instances of a general class of failures
arising when complex systems operate under insufficient constraint. The formal machinery
introduced encompasses projection-based inference over high-dimensional configuration spaces,
prior-dominant attractor dynamics, entropy-driven diffusion following constraint boundary
violations, and the condition of constraint closure as the determinant of whether a system
produces a meaningful output or degenerates into an arbitrary sequence of locally admissible
steps. The analysis demonstrates that such failures are not isolated accidents but structurally
predictable consequences of partial constraint in rapidly iterating, high-dimensional systems.
The essay concludes by deriving parallel design implications for software pipelines and
argumentative practice, and by situating both within a general principle governing the
reliability of observable artifacts as evidence of underlying generative structure.
\end{abstract}
\end{titlepage}

\newpage
\tableofcontents
\newpage

% ==========================================================
\section{Introduction}
% ==========================================================

This essay examines two related phenomena through a unified formal framework. The first is
a recent incident in which proprietary AI system code was accidentally released into the
public domain, an event widely interpreted as a failure of operational discipline but here
analyzed as a structurally predictable outcome of high-dimensional system dynamics. The second
is the degeneration of essay form under generative pressure, a phenomenon in which the surface
properties of argumentative writing are reproduced without the underlying constraint structure
that gives those properties meaning.

The central claim is that both phenomena are instances of a general class of failures arising
when complex systems operate under insufficient constraint. The formal machinery developed to
analyze the code leak---encompassing projection-based inference, prior-dominant attractor
dynamics, entropy-driven diffusion, and the condition of constraint closure---applies without
modification to the epistemological analysis of written discourse. This unification is not
merely analogical. It reflects a deeper structural identity between the dynamics of software
pipelines and the dynamics of argumentative trajectories.

In both domains, the observable surface exceeds its underlying necessity. A running system and
a grammatically fluent essay are both low-dimensional projections of the processes that
produced them. The reliability of either as evidence of the internal structure it represents
depends on whether the generating process was sufficiently constrained to permit
near-identification of that structure from the observable output. When constraint is partial,
surfaces decouple from interiors. The result, in the case of software, is a pipeline that
admits trajectory escape into unintended configurations; in the case of writing, a sequence
of sentences that resembles argument without instantiating it.

The essay proceeds as follows. Sections 2 through 7 develop the formal analysis of the code
leak, introducing the projection framework, the near-identifiability result, attractor
dynamics, entropy diffusion, constraint failure in dynamical workflows, and the economics of
computational enclosure. Section 8 establishes the bridging argument, formalizing essay
writing as a constrained dynamical process and making explicit the structural identity that
motivates the unified treatment. Sections 9 through 14 apply the formal machinery to the
epistemological domain, analyzing the degeneration of essay form, the structural basis of
clarity, the collapse of length as a reliability signal, and the shift in trust from artifact
to process. Section 15 addresses the discourse on safety and the error of reasoning across
incompatible abstraction levels. Section 16 derives design implications. Section 17 concludes.

% ==========================================================
\section{The Incident as a Projection Event}
% ==========================================================

Let $\config$ denote the full internal configuration space of an AI development system,
understood as the product of all relevant state dimensions: source code, build artifacts,
workflow specifications, deployment scripts, infrastructure definitions, organizational
protocols, and the latent practices that govern their composition. A point $X \in \config$
represents a complete instantaneous specification of the system.

No external observer has access to $X$ directly. What is observable is a projection
$\proj(X) \in \obs$, where $\obs$ is a lower-dimensional space and $\proj : \config \to \obs$
is the projection map induced by the particular leak event. In the incident under analysis,
$\proj(X)$ consists of the released code artifacts together with whatever structural
inferences can be drawn from their organization, naming conventions, and dependency structure.

\begin{definition}[Observation-consistent set]
Given an observation $o \in \obs$, the observation-consistent set is
\[
  \config_o = \{ X \in \config : \proj(X) = o \}.
\]
This is the fiber of $\proj$ over $o$, and represents all internal configurations that would
have produced the observed projection.
\end{definition}

The central interpretive question is whether $\config_o$ is sufficiently constrained to
permit reconstruction of the underlying system. In general, $\proj$ fails to be injective:
multiple configurations produce the same projection, and $\config_o$ may be high-dimensional.
What the incident reveals, however, is not the impossibility of reconstruction but rather
the regime of near-identifiability, in which the fiber is structured enough to admit
efficient convergence to configurations that are functionally indistinguishable from $X$.

\begin{proposition}[Near-identifiability under architectural priors]
Let $\prior$ denote the prior distribution over $\config$ induced by shared architectural
conventions, published design patterns, and existing open-source systems. If the support of
$\prior$ restricted to $\config_o$ concentrates around a small number of equivalence classes
under functional equivalence, then practical reconstructibility holds even when strict
injectivity of $\proj$ fails.
\end{proposition}

The rapid emergence of working replicas following the incident constitutes empirical evidence
for this proposition. Reconstruction was not performed by exhaustive search over $\config_o$
but by gradient descent in the prior landscape, converging to the nearest attractor basin
consistent with the observed projection.

% ==========================================================
\section{Projection Degeneracy and Near-Identifiability}
% ==========================================================

The fiber $\config_o$ is in general a manifold of dimension $\dim \config - \dim \obs$.
When $\config$ is high-dimensional and $\proj$ collapses many dimensions, this fiber may
itself be high-dimensional. The degeneracy of the projection is thus measured by the
codimension of $\obs$ in $\config$.

In the present case, the projection exposes a significant but still partial slice of $X$.
The released artifacts include module boundaries, interface signatures, and dependency
relationships sufficient to constrain the space of admissible implementations. However, many
internal decisions remain hidden: the specific optimization choices, the training
infrastructure, the evaluation pipeline, and the coordination protocols between subsystems.
The fiber $\config_o$ is therefore high-dimensional, but not uniformly so.

\begin{definition}[Effective degeneracy]
Let $\config_o^{\prior}$ denote the restriction of the prior $\prior$ to $\config_o$.
The effective degeneracy of the projection is the entropy
\[
  H(\config_o^{\prior}) = -\int_{\config_o} \prior(X) \log \prior(X) \, dX.
\]
Low effective degeneracy means the prior concentrates the fiber onto a small number of
typical configurations, enabling reconstruction even under high nominal degeneracy.
\end{definition}

The distinction between nominal degeneracy---the dimension of the fiber---and effective
degeneracy---the entropy of the prior-weighted fiber---is central to the analysis. A
projection may be nominally degenerate while admitting near-identification because the
prior structure of the domain strongly constrains which configurations are typical.

This result has a general implication: the security value of enclosure does not scale with
the nominal dimension of the hidden configuration space, but with the effective degeneracy
of the projection under the shared prior of the reconstruction domain. In a mature field
with strong architectural conventions, this effective degeneracy may be surprisingly low
even when the nominal configuration space is vast.

% ==========================================================
\section{Prior-Dominant Attractors and Rapid Reconstruction}
% ==========================================================

The reconstruction of a system from its projection is most naturally understood as a
dynamical process on $\config$. External agents, operating within a shared domain
characterized by learned architectural priors, initialize their reconstruction at some
point $X_0$ and evolve according to an energy functional that balances fidelity to the
observed projection against consistency with prior expectations.

\begin{definition}[Reconstruction energy functional]
Let $\lambda > 0$ be a regularization parameter. The reconstruction energy functional is
\[
  \energy(X) = \| \proj(X) - o \|^2 + \lambda \cdot (-\log \prior(X)),
\]
where the first term penalizes deviation from the observed projection and the second term
penalizes departure from the prior. Reconstruction trajectories evolve by gradient descent
on $\energy$.
\end{definition}

The attractor structure of $\energy$ is determined primarily by the prior term when $\lambda$
is large relative to the constraint strength of the projection. In this prior-dominant
regime, the landscape of $\energy$ is essentially the landscape of $-\log \prior$, and
trajectories converge to the modes of the prior that are consistent with the projection.

\begin{definition}[Prior-dominant attractor]
A configuration $X^* \in \config_o$ is a prior-dominant attractor if it is a local minimum
of $\energy$ in the prior-dominant regime, i.e., if $X^*$ is a mode of $\prior$ restricted
to $\config_o$.
\end{definition}

The rapid emergence of working replicas following the incident is explained by this
structure. The reconstruction problem is not one of deriving the hidden system from first
principles but of identifying the prior-dominant attractor in the neighborhood of the
observed projection. When the prior is sufficiently concentrated---as it is in a mature
field with strong shared conventions---this identification is computationally efficient.

\begin{theorem}[Convergence under prior dominance]
Suppose $\prior$ is log-concave and the projection $\proj$ is linear. Then the reconstruction
energy $\energy$ is convex on $\config_o$ in the prior-dominant regime, and gradient descent
converges to the unique global minimum at a rate determined by the condition number of the
Hessian of $\energy$ restricted to $\config_o$.
\end{theorem}

\begin{proof}
Under log-concavity, $-\log \prior$ is convex. The quadratic fidelity term is convex by
construction. The restriction of a convex function to a linear fiber is convex. Convergence
of gradient descent on a convex function follows from standard results, with rate determined
by the ratio of the largest to smallest eigenvalues of the Hessian.
\end{proof}

The practical import of this result is that near-identifiability and rapid reconstruction
are not properties of the leaked artifact alone, but of the interaction between the artifact
and the prior landscape of the reconstruction domain. Enclosure is therefore structurally
fragile in any domain where shared priors are sufficiently expressive.

% ==========================================================
\section{Entropy Diffusion and Structural Relaxation}
% ==========================================================

Once the constraint boundary of the enclosure has been breached and the projection $\proj(X)$
has entered the observable domain, the subsequent dissemination of the leaked material
follows the dynamics of entropy-driven diffusion on a networked information environment.

Let $\rho(n, t)$ denote the density of the leaked information at node $n$ at time $t$ in
the network. In the continuous approximation, the evolution of $\rho$ is governed by the
diffusion equation
\[
  \frac{\partial \rho}{\partial t} = D \nabla^2 \rho - \mu \rho,
\]
where $D$ is the diffusion coefficient determined by network connectivity and sharing
propensity, and $\mu$ is a decay term accounting for information becoming stale or
superseded. In the absence of active counter-dissemination, the system evolves toward a
diffuse equilibrium in which the leaked material has propagated throughout the reachable
component of the network.

This process is irreversible in the thermodynamic sense. The total entropy of the
information distribution,
\[
  S(t) = -\int \rho(n,t) \log \rho(n,t) \, dn,
\]
is monotonically non-decreasing under the diffusion dynamics. Once the constraint boundary
is breached, the system cannot return to the low-entropy state in which the information was
concentrated within the enclosure. Retraction of the leaked material is not a reversal of
the diffusion process but an attempt to impose a new boundary condition on an already
relaxed distribution, an intervention that cannot undo the structural relaxation that has
occurred.

\begin{proposition}[Irreversibility of entropic relaxation]
Let $\rho_0$ be the initial distribution concentrated within the enclosure, and let $\rho_t$
be the distribution at time $t > 0$ under diffusion dynamics. Then $S(\rho_t) > S(\rho_0)$
for all $t > 0$, and there is no admissible dynamics on the network that returns the system
to $\rho_0$ without external constraint reimposition.
\end{proposition}

The general principle that emerges is one of entropic smoothing: localized constraint
failures result in global redistribution of informational content until a new equilibrium
is reached. The significance of the incident lies not in the quantum of information released
but in the irreversibility of the relaxation it initiates.

% ==========================================================
\section{Constraint Failure in Dynamical Workflows}
% ==========================================================

The conditions under which the incident occurred can be formalized by treating the
deployment pipeline as a constrained dynamical system. Let the state of the pipeline at
time $t$ be a point $s_t$ in a state space $S$. The admissible states form a submanifold
$S_{\text{adm}} \subset S$ defined by the constraints imposed by security protocols, access
controls, and workflow specifications. The pipeline is designed to evolve along trajectories
$\traj : [0, T] \to S$ that remain within $S_{\text{adm}}$.

\begin{definition}[Constraint closure]
A pipeline achieves constraint closure if for every admissible initial state $s_0 \in S_{\text{adm}}$
and every admissible input sequence, the resulting trajectory $\traj$ satisfies
$\traj(t) \in S_{\text{adm}}$ for all $t \in [0, T]$.
\end{definition}

The presence of manual steps in the pipeline introduces unresolved degrees of freedom.
A manual step is one whose output is not determined by the preceding state and a fixed
computational rule, but depends on the judgment, attention, and situational awareness of
a human operator. Such steps prevent full constraint closure because the operator's action
cannot be guaranteed to preserve admissibility under all conditions.

\begin{proposition}[Trajectory escape under partial closure]
Let a pipeline contain $k$ manual steps. The probability that the trajectory remains within
$S_{\text{adm}}$ over the full execution is at most $\prod_{i=1}^k p_i$, where $p_i$ is the
probability that the operator at step $i$ takes an admissible action. If any $p_i < 1$,
the pipeline does not achieve constraint closure, and trajectory escape occurs with positive
probability.
\end{proposition}

The incident is therefore a trajectory-level failure rather than a static error. The system
had not achieved structural completion: the configuration was not the result of a constraint
violation in a closed system but of an admissibility condition that had never been fully
enforced. The distinction is important because it changes the appropriate remediation
strategy. Static errors call for correction; trajectory-level failures call for redesign
of the constraint structure.

% ==========================================================
\section{The Enclosure of Computation and Its Breakdown}
% ==========================================================

The economic significance of the incident is best understood through the concept of
computational enclosure. Proprietary AI systems derive their commercial value not from the
intrinsic scarcity of the computational processes they implement---since these processes
can, in principle, be replicated by any agent with sufficient prior knowledge and
computational resources---but from the enclosure of access to those processes within a
controlled interface.

Computational enclosure functions analogously to other forms of enclosure in economic
history: it converts a process that would otherwise be freely reproducible into a commodity
whose access can be regulated and monetized. The value of the enclosure is a function of
the cost of reconstruction from outside, which in turn depends on the effective degeneracy
of the projection available to external agents.

The near-identifiability result from Section 3 implies that this enclosure is structurally
fragile in any domain where the shared prior of potential reconstructors is sufficiently
expressive. As the prior landscape of the AI development domain becomes richer---through
the accumulation of published research, open-source systems, and trained practitioners---the
marginal value of the constraint boundary decreases. The enclosure depends for its
effectiveness on the reconstruction cost remaining high, but reconstruction cost is a
decreasing function of prior expressiveness.

The leak temporarily collapses this enclosure, enabling widespread replication and
recomposition of the leaked system's functional properties. The economic consequence is not
primarily the loss of the specific information contained in the projection but the
demonstration that the enclosure is weaker than its market valuation implied. This
informational effect on beliefs about future reconstruction costs may be more significant
than the direct informational content of the leak itself.

% ==========================================================
\section{Degeneracy, Recoverability, and the Collapse of Enclosure}
% ==========================================================

The preceding analysis introduces three concepts that, while developed in different
sections, describe a single underlying phenomenon: effective degeneracy of projection,
practical recoverability of the generating system, and the stability of computational
enclosure. This section establishes their formal relationship.

Let $o = \proj(X)$ be an observed projection of a configuration $X \in \config$, and
let $\config_o^{\prior}$ denote the prior-weighted fiber over $o$. Recall that the
effective degeneracy of the projection is given by
\[
  H(\config_o^{\prior}).
\]

\begin{definition}[Reconstruction cost]
Let $\mathcal{A}$ be a class of reconstruction algorithms operating under prior
$\prior$. The reconstruction cost $C(o)$ is the expected computational effort required
for $\mathcal{A}$ to produce a configuration $\hat{X}$ such that
\[
  \proj(\hat{X}) = o \quad \text{and} \quad \hat{X} \sim X \text{ under functional equivalence}.
\]
\end{definition}

\begin{proposition}[Degeneracy--cost relation]
Under a fixed prior $\prior$ and reconstruction class $\mathcal{A}$, the reconstruction
cost $C(o)$ is a monotone increasing function of the effective degeneracy
$H(\config_o^{\prior})$.
\end{proposition}

\begin{proof}
Low effective degeneracy implies that $\config_o^{\prior}$ is concentrated on a small
number of high-probability configurations. Gradient-based or sampling-based procedures
converge rapidly to these modes. As degeneracy increases, the mass of $\config_o^{\prior}$
spreads over a larger set, increasing the search complexity required to locate a
functionally equivalent configuration.
\end{proof}

This establishes that effective degeneracy is the fundamental quantity governing
recoverability: systems with low degeneracy are practically reconstructible even when
the projection is not injective.

\begin{definition}[Enclosure strength]
The strength of a computational enclosure is defined as
\[
  E = C(o),
\]
the reconstruction cost of the system given its observable projection.
\end{definition}

This definition captures the economic function of enclosure: it derives its value
from the difficulty of external reconstruction rather than from intrinsic properties
of the system.

\begin{theorem}[Degeneracy--enclosure equivalence]
Under fixed prior $\prior$, the following are equivalent up to monotone transformation:
\begin{enumerate}[leftmargin=2em]
\item Low effective degeneracy $H(\config_o^{\prior})$
\item Low reconstruction cost $C(o)$
\item Weak enclosure strength $E$
\end{enumerate}
\end{theorem}

\begin{proof}
The equivalence of (1) and (2) follows from the degeneracy--cost relation. The
equivalence of (2) and (3) is immediate from the definition of enclosure strength.
Monotonicity preserves ordering across all three quantities.
\end{proof}

This result formalizes the collapse of enclosure observed in the incident. The leak
did not merely expose information; it revealed that the effective degeneracy of the
projection under the domain prior was already low. Reconstruction was therefore
inevitable, and the enclosure was structurally weak prior to the leak.

\subsection{Prior Dependence and Temporal Instability}

The equivalence above holds relative to a fixed prior $\prior$. In practice, $\prior$
evolves over time as the domain accumulates knowledge.

\begin{proposition}[Prior-induced collapse]
Let $\prior_t$ be the prior at time $t$, with increasing expressiveness over time.
Then for a fixed projection $o$,
\[
  H(\config_o^{\prior_t}) \downarrow \quad \Rightarrow \quad C_t(o) \downarrow \quad \Rightarrow \quad E_t \downarrow.
\]
\end{proposition}

Thus, even in the absence of new leaks, enclosure strength decays over time as the
shared prior of the field becomes richer.

\subsection{Interpretation}

The collapse of computational enclosure is therefore not an event but a process. The
leak acts as a discontinuity that reveals a pre-existing condition: that the system
resides in a regime of low effective degeneracy under contemporary priors.

The same structure applies to argumentative texts. A text with low constraint density
has high effective degeneracy, making its generating process unrecoverable and its
content informationally weak. A text with high constraint density has low degeneracy,
making its generating process identifiable and its structure resistant to arbitrary
reconstruction.

In both cases, the governing principle is identical: the reliability and value of a
projection are determined not by its surface complexity, but by the extent to which
it reduces the space of admissible generating processes.

% ==========================================================
\section{Speed, Automation, and the Limits of Control}
% ==========================================================

The incident can be further situated within the general dynamics of high-velocity development
environments. Such environments are characterized by the combination of partial automation
with human-in-the-loop processes, a combination that creates a regime in which constraint
enforcement is probabilistic rather than absolute.

Let $v$ denote the iteration velocity of the development pipeline, measured by the number
of state transitions per unit time. As $v$ increases, the time available for human
operators to verify admissibility at each manual step decreases proportionally. If the
cognitive cost of verification is approximately constant per step, and the time budget per
step is $1/v$, then the probability of successful verification at each step is a decreasing
function of $v$. Combining this with the result of Section 6:

\begin{corollary}[Constraint enforcement probability under velocity pressure]
In a pipeline with $k$ manual steps operating at velocity $v$, the probability that the
trajectory remains within $S_{\text{adm}}$ decreases monotonically with $v$. For sufficiently
high $v$, trajectory escape becomes effectively certain over a long enough operational period.
\end{corollary}

This result formalizes the intuition that such failures are emergent properties of the
system architecture rather than isolated lapses of individual operators. The operator at
the relevant step did not fail to apply a rule; the system was operating in a regime in
which reliable rule application was structurally impossible given the time constraints
imposed by the development velocity. The appropriate unit of analysis is therefore not the
individual action but the dynamical regime.

The tension between iteration speed and reliability is not resolvable within the existing
architecture: increasing speed while maintaining manual steps produces a system in which
constraint enforcement probability goes to zero. Resolution requires either reducing velocity,
eliminating manual steps through automation, or restructuring the constraint system so that
admissibility is enforced by structural invariants rather than by operator verification.

% ==========================================================
\section{Writing as a Constrained Dynamical Process}
% ==========================================================

The two domains examined in this essay---the structural analysis of AI system failures and
the epistemological critique of essay form under generative pressure---are not merely
analogous. They are instances of the same underlying phenomenon: the behavior of
high-dimensional systems under partial constraint. Establishing this identity requires
showing that the formal machinery developed in the preceding sections applies without
modification to the dynamics of argumentative writing.

An essay and a software pipeline share a common formal structure. Both are trajectories
through a space of possible states, where the admissibility of each transition depends on
conditions imposed by prior commitments. In a software pipeline, these conditions are
encoded in interfaces, contracts, and workflow dependencies. In an argumentative text, they
are encoded in premises, definitions, and logical entailments. In both cases, constraint
closure---the condition in which all degrees of freedom are resolved by the governing
structure---determines whether the system produces a meaningful output or degenerates into
an arbitrary sequence of locally plausible steps.

The failure modes are structurally parallel. A pipeline with unresolved manual steps admits
transitions that violate admissibility, producing trajectories that escape the intended
configuration space. An essay produced without governing premises admits continuations that
are locally coherent but globally arbitrary, producing text that resembles argument without
instantiating it. In both cases, the observable surface, whether running code or grammatical
prose, provides insufficient evidence of the underlying constraint structure.

The projection framework introduced in Section 2 applies directly to textual artifacts.
A reader observing a finished essay sees only a low-dimensional projection of the
argumentative process that produced it. The reconstructibility of that process---and therefore
the trustworthiness of the artifact---depends on whether the projection is sufficiently
constrained to permit near-identification of the underlying trajectory. Transparent reasoning,
explicit premises, and traceable inference steps are precisely the structural features that
reduce projection degeneracy in written discourse, concentrating the observation-consistent
set around a small number of admissible argumentative trajectories.

Similarly, the attractor dynamics of Section 4 apply to the formation of arguments under
generative pressure. Shared conceptual priors, disciplinary conventions, and rhetorical
patterns constitute an attractor landscape within which generated text settles. Without
external constraint, generation converges to the nearest basin regardless of whether that
basin corresponds to a true argument. The result is the prior-dominant essay: structurally
familiar, disciplinarily conventional, intellectually inert.

\begin{definition}[Argumentative constraint closure]
An essay achieves argumentative constraint closure if for every premise in the governing
structure, all transitions between claims are determined by the logical and evidential
relations specified by that structure, and no claim appears whose inclusion or position
is underdetermined by the constraints.
\end{definition}

Without constraint closure, the essay is formally indistinguishable from an unconstrained
generation: locally plausible, globally arbitrary, and incapable of serving as reliable
evidence of the reasoning process that produced it.

\section{The Necessity of Prompted Structure}

The analysis of writing as a constrained dynamical process has an immediate operational
implication: the constraint structure of an essay cannot be inferred post hoc from an
unconstrained generative request. It must be specified prior to or during the generation
process. A request of the form ``write an essay on topic $T$'' defines only a projection
target in $\obs$ and leaves the corresponding fiber in $\config$ maximally degenerate.

\begin{definition}[Unconstrained essay prompt]
A prompt is unconstrained if it specifies a target topic or domain without imposing
conditions on the admissible argumentative trajectories, such as required premises,
explicit tensions, or exclusion criteria.
\end{definition}

Under an unconstrained prompt, the generating process operates entirely within the prior
landscape described in Section 4. The resulting text is therefore a realization of a
prior-dominant attractor rather than a trajectory determined by a specific constraint
structure. The observable output may satisfy the formal properties of an essay, but the
underlying trajectory is not identifiable from the projection, since the same output class
is consistent with a large set of generating processes.

By contrast, a constrained prompt specifies a submanifold of admissible trajectories in
$\config$. This reduces the effective degeneracy of the projection and forces the generator
to produce a trajectory that is determined, at least partially, by the imposed constraints.

\begin{proposition}[Constraint necessity for argumentative generation]
Let $\mathcal{G}$ be a generative system operating over $\config$. Then $\mathcal{G}$
produces an argumentatively identifiable trajectory if and only if the prompt defines a
constraint set whose projection reduces the effective degeneracy of the observation-consistent
set.
\end{proposition}

The implication is that the responsibility for producing an essay does not lie solely with
the generator but with the specification of the constraint structure. Without such
specification, the generator cannot produce necessity, because necessity is not a property
of the output distribution but of the constraint space within which the output is selected.

% ==========================================================
\section{The Degeneration of Essay Form Under Unconstrained Generation}
% ==========================================================

The contemporary availability of generative systems has introduced a fundamental asymmetry
between form and substance in written discourse. It is now trivial to produce text that
satisfies the superficial criteria of an essay: paragraph organization, grammatical
continuity, disciplinary vocabulary, and transitions that simulate logical progression.
These properties no longer reliably indicate the presence of an underlying argument,
because they can be produced without one.

An essay, properly understood, is not defined by its length or its surface organization
but by the existence of a constrained argumentative trajectory. Each claim is positioned
relative to others in such a way that the structure is not arbitrary: the removal or
reordering of components alters the meaning of the whole. This internal dependency is what
distinguishes argument from aggregation.

Unconstrained generation severs this dependency. When text is produced without a governing
structure of premises, definitions, or tensions, it becomes interchangeable in the sense
relevant to the projection framework: the observation-consistent set over the space of
generating processes is maximally degenerate, since the surface output is consistent with
the entire range of unconstrained generative processes. The result is an essay-shaped
object that lacks necessity. Its sentences follow one another, but not because they must.

This is experienced by readers as a form of cognitive emptiness, despite the appearance of
completeness. The experience is not aesthetic but epistemic: the reader correctly perceives
that the text fails to reduce uncertainty about the argumentative trajectory that produced
it. No information about the generator is available from the output, because the output is
consistent with all generators. The text is maximally uninformative in the
information-theoretic sense while appearing maximally structured in the surface sense.

% ==========================================================
\section{Constraint as the Source of Clarity}
% ==========================================================

The widespread claim that contemporary audiences prefer brevity over long-form writing is
frequently misinterpreted as a rejection of essays as a form. A more precise characterization
is that readers are rejecting unconstrained verbosity. What is selected for in the current
environment is not shortness but clarity, and clarity is a structural rather than aesthetic
property. It emerges when the degrees of freedom in an argument have been sufficiently
reduced by the governing constraint structure.

\begin{definition}[Structural clarity]
The structural clarity of an argumentative text is inversely proportional to the effective
degeneracy of its projection over the space of generating processes. A text is structurally
clear to the degree that its observable properties constrain the space of argumentative
trajectories consistent with it.
\end{definition}

Under this definition, a short text can be structurally opaque if its brevity is achieved
by compression that eliminates the inferential structure that would permit reconstruction,
and a long text can be structurally clear if each sentence functions as an admissibility
constraint on the subsequent argument. The relevant quantity is not length but constraint
density: the ratio of constraint-imposing content to total content.

Constraint reduces the space of admissible statements locally and globally. When an argument
is governed by explicit premises and a defined problem, each sentence must contribute to
resolving the tension established by that structure. Irrelevant or redundant expressions are
excluded not by editorial judgment but by the logic of the constraint space: they are
inadmissible transitions. The resulting text appears clearer because it is necessary, and
it is trustworthy because its observable properties reduce the degeneracy of its projection
over generating processes.

% ==========================================================
\section{Constraint Density as a Design Variable}
% ==========================================================

The preceding sections treat constraint density as an emergent property of a
completed argumentative artifact. For practical purposes, however, constraint
density must be understood as a design variable: a quantity that can be actively
controlled during the generation process.

Let $Y$ be a text generated under a constraint specification $C$. Then the
resulting constraint density $D(Y)$ is a function of both the richness of $C$
and the fidelity with which the generating process adheres to it. Formally,
\[
  D(Y) = D(C, \mathcal{G}),
\]
where $\mathcal{G}$ denotes the generator.

A weak constraint specification---one that defines a topic but not a trajectory---induces
a high-entropy trajectory space $\mathcal{T}(Y)$, resulting in low constraint density.
A strong specification---one that defines premises, tensions, and admissibility
conditions---reduces $\mathcal{T}(Y)$ and increases $D(Y)$.

This reframes writing as a problem of constraint design rather than content production.
The relevant question is not what sentences should appear, but what degrees of freedom
must be eliminated such that only a narrow class of trajectories remains admissible.

\begin{proposition}[Constraint density monotonicity]
Let $C_1 \subset C_2$ be two constraint sets, where $C_2$ imposes strictly more
restrictions on admissible trajectories than $C_1$. Then for any generator $\mathcal{G}$,
\[
  D(Y_{C_2}) \geq D(Y_{C_1}),
\]
where $Y_{C_i}$ is the text generated under constraint set $C_i$.
\end{proposition}

This result implies that increasing clarity is not a matter of refining expression
but of eliminating admissible alternatives. Clarity is therefore achieved not by
adding content, but by reducing possibility.

In practical terms, the design of an essay becomes the design of a constraint surface
within the trajectory space. The text that results is a trace of that surface, and
its clarity is determined by how sharply the surface restricts movement.

% ==========================================================
\section{Essay as Trajectory Rather Than Artifact}
% ==========================================================

The foregoing analysis supports a reconceptualization of the essay as a trajectory through
a constrained space rather than as a static artifact. The value of an essay lies in the
path it traces: how it moves from initial assumptions to conclusions, how it navigates and
excludes alternative possibilities, and how each transition is determined by the constraint
structure rather than by the generator's local preferences.

\begin{proposition}[Necessity as the criterion of argumentative value]
An argumentative text has positive necessity if and only if there exists at least one
admissible alternative at each transition that is excluded by the governing constraint
structure. A text of zero necessity is one in which all transitions are equally admissible,
and the specific trajectory traced is determined only by the prior distribution of the
generator.
\end{proposition}

This proposition makes precise the intuition that a real argument has directionality: if
a step is removed, something breaks; if the order is reversed, coherence degrades. This
fragility is a signal of rigor rather than a defect, because it indicates that the
trajectory is constrained. A text that could be rearranged without consequence is one in
which the constraint structure is empty.

The implication for practice is that writing cannot be separated from the specification of
constraints. To request an essay without providing an argument is to request a trajectory
without defining a constraint space. The resulting text may exhibit local coherence---each
sentence is locally plausible---but lacks global direction. The formal parallel with
pipeline constraint failure is exact: in both cases, the absence of closure admits arbitrary
continuation, and the resulting trajectory may escape the intended configuration space
entirely without any single step being individually inadmissible.

\section{Perceived Emptiness and the Failure of Informational Gain}

The reader's experience of unconstrained generative text as ``empty'' can be formalized in
information-theoretic terms. Let $Y$ denote the observable text and $G$ the generating
process. The informational value of $Y$ with respect to $G$ is given by the mutual information
$I(Y; G)$.

In the case of a constrained argumentative trajectory, $Y$ reduces uncertainty about $G$,
because the constraint structure leaves a detectable imprint on the observable output.
Different generating processes would produce observably different trajectories under the
same constraints. Thus $I(Y; G)$ is positive.

In the case of unconstrained generation, however, the observable output is consistent with
a large class of generating processes. The same essay-shaped text could have been produced
by many different trajectories within the prior landscape. The mutual information
$I(Y; G)$ is therefore low, approaching zero in the limit of maximal degeneracy.

\begin{definition}[Informational emptiness]
An argumentative text is informationally empty if $I(Y; G)$ is below a threshold required
for reliable inference of the generating process.
\end{definition}

This formalization explains why readers disengage from such text. The act of reading is an
attempt to reconstruct the generating trajectory from the observable projection. When the
projection is maximally degenerate, this reconstruction fails, and the reader receives no
informational gain about the underlying reasoning process.

The preference for shorter, distilled content is therefore not a preference for reduced
content but for increased information density. A shorter text that imposes strong
constraints may yield higher $I(Y; G)$ than a longer unconstrained text. The relevant
quantity is not length but the reduction in uncertainty about the generating trajectory
achieved by observing the text.

% ==========================================================
\section{The Collapse of Length as a Signal}
% ==========================================================

Historically, the length and structural organization of a text functioned as low-dimensional
proxies for a high-dimensional property: the intellectual effort, sustained engagement, and
constraint-navigating work required to produce a genuine argument. The reliability of these
proxies depended on the production costs of long structured text being sufficiently high to
make simulation expensive. Under these conditions, form and substance were correlated in
expectation, and readers could rationally treat surface features as evidence of underlying
argumentative quality.

Generative systems have broken this correlation by enabling high-dimensional surface
properties to be produced from a near-zero-constraint process. This is precisely the
situation analyzed in Section 3 for the code leak: the projection has been decoupled from
the configuration it was formerly reliable evidence of. Length and structural organization
now belong to the class of low-cost observables that can be produced without instantiating
the high-dimensional property they previously tracked.

The reader's response to this decoupling is structurally rational. When a proxy variable
loses its correlation with the quantity it proxied, agents who relied on the proxy update
their inference strategy. Readers are increasingly attending to properties that remain
difficult to simulate under unconstrained generation: internal consistency over long
dependencies, explicit acknowledgment of counter-arguments, traceable inferential structure,
and the visibility of constraint. These properties are more resistant to unconstrained
generation because they impose structural requirements on the generating process itself,
not merely on the observable output.

The collapse of length as a signal is therefore not a cultural shift in attention spans
but a rational Bayesian update in response to a change in the generating environment. The
prior over generators, given an observation of long structured text, has shifted: whereas
previously it was concentrated on processes involving sustained argumentative constraint,
it is now distributed across a much larger class of processes, most of which involve no
constraint at all.

% ==========================================================
\section{Adversarial Generation and the Simulation of Constraint}
% ==========================================================

A potential objection to the preceding analysis is that generative systems are not
limited to unconstrained outputs. They can simulate many of the surface properties
associated with constraint: explicit premises, structured arguments, and even
counter-argumentation. This raises the question of whether constraint density can
itself be simulated without being instantiated.

Let $Y$ be a text that exhibits surface markers of constraint. The question is
whether these markers correspond to actual reductions in the admissible trajectory
space or merely to patterns that resemble such reductions.

\begin{definition}[Simulated constraint]
A text exhibits simulated constraint if it contains surface features associated
with constraint (e.g., explicit premises or logical transitions) without those
features reducing the effective degeneracy of the observation-consistent set.
\end{definition}

Simulated constraint arises when the generator produces patterns that are correlated
with constraint under the training distribution but does not enforce the underlying
relations that give those patterns necessity. In such cases, the apparent structure
does not restrict the trajectory space, because alternative continuations remain
equally admissible.

\begin{proposition}[Indistinguishability under shallow observation]
If observation is restricted to local or short-range dependencies, then simulated
constraint is observationally indistinguishable from genuine constraint.
\end{proposition}

\begin{proof}
Local patterns associated with constraint can be reproduced independently of the
global structure. If observation does not capture long-range dependencies or
global consistency, then both genuine and simulated constraint produce identical
local statistics.
\end{proof}

The distinction becomes observable only under deeper inspection: long-range
dependency tracking, counterfactual perturbation (removing or altering a premise),
or attempts to reconstruct the generating trajectory. Genuine constraint produces
fragility under such perturbations, while simulated constraint remains robust
because no underlying dependency exists.

This introduces an adversarial dimension to the evaluation of argumentative text.
As generators improve, the detection of constraint requires increasingly global
tests, shifting the burden from surface evaluation to structural analysis.

% ==========================================================
\section{Transparency, Bias, and Trust in High-Volume Text Environments}
% ==========================================================

In an environment saturated with generated text, the trust function of a reader cannot be
adequately defined over the artifact alone. The reader must attempt to infer the constraint
structure of the generating process from the available evidence, and the reliability of this
inference depends on the degree to which the artifact makes that structure legible.

\begin{definition}[Process legibility]
An argumentative text has high process legibility if its observable properties substantially
reduce the effective degeneracy of its projection over the space of generating processes.
Equivalently, high legibility means that the observation-consistent set is concentrated
around a small number of generating trajectories with similar constraint structures.
\end{definition}

Transparency in this context is the act of increasing process legibility by making the
constraint structure of the argument explicit. An argument that exposes its premises,
acknowledges its exclusions, enumerates the alternative positions it considered, and traces
its inferential steps reduces projection degeneracy by providing observational content that
is inconsistent with most unconstrained generating processes. The underlying trajectory
becomes more nearly identifiable.

Transparency does not eliminate bias in the sense of eliminating the generator's particular
perspective or prior commitments. It makes bias legible by making the constraint structure
explicit. A reader who can identify the constraints under which an argument was produced
can evaluate its conclusions relative to those constraints, rather than treating the argument
as if it were constraint-free. This is the appropriate epistemic stance in a high-volume
generative environment: not the demand for unbiased text, which is formally incoherent, but
the demand for text whose biases are identifiable from its observable properties.

The preference, widely observed in contemporary information practice, for distilled summaries
from sources with known perspectives over comprehensive neutral accounts from anonymous
generators, is therefore not a retreat from epistemic standards but an adaptation to a
changed environment. When the volume of available text exceeds the capacity for detailed
process reconstruction, readers rely on reputational signals as proxies for constraint
structure. The source's known perspective functions as a prior over the constraint structures
they typically employ, reducing the effective degeneracy of the projection without requiring
the reader to verify the constraint structure from the text alone.

\section{Reputation as a Prior Over Constraint Structures}

In high-volume text environments, readers cannot reconstruct the generating process for
each observed artifact. Instead, they rely on priors over classes of generators. A source's
reputation functions as such a prior, encoding expectations about the constraint structures
that the source typically employs.

\begin{definition}[Reputational prior]
Let $\mathcal{G}$ be the space of generating processes. A reputational prior is a probability
distribution $\pi$ over $\mathcal{G}$ conditioned on the identity of the source, reflecting
the expected constraint structures of that source's outputs.
\end{definition}

When a reader encounters a text from a known source, the effective degeneracy of the
projection is reduced by conditioning on $\pi$. The observation-consistent set becomes
\[
  \config_{Y}^{\pi} = \{ G \in \mathcal{G} : G \text{ could have produced } Y \text{ and is
  consistent with } \pi \}.
\]

If $\pi$ is sufficiently concentrated, then even a relatively compressed or partial
artifact can yield high process legibility. Conversely, text from an unknown or
unconstrained generator corresponds to a diffuse prior, leaving the projection highly
degenerate and reducing trust.

This explains the empirical preference for distilled summaries from sources with known
biases over longer texts from anonymous or opaque generators. Bias, when made explicit,
functions as a constraint on the generating process, thereby reducing degeneracy. The
reader can condition on this constraint and evaluate the argument within its proper frame.

The absence of declared bias does not eliminate bias but increases degeneracy by removing
a constraint that could otherwise be used to infer the generating trajectory. In this
sense, transparency is not a moral property but an informational one: it reduces the
uncertainty of the inverse problem faced by the reader.

% ==========================================================
\section{Misplaced Abstraction and the Discourse on Safety}
% ==========================================================

Public interpretations of both the code leak and the degeneration of essay form exhibit a
characteristic error of reasoning: the conflation of observations at one abstraction level
with conclusions that are only warranted at another. In the case of the leak, local
operational failures are treated as evidence for global claims about the safety properties
of AI systems. In the case of essay form, the local observation that generative systems
can produce essay-shaped text is treated as evidence for global claims about the nature of
reasoning and the death of argumentative culture.

Both errors have the same formal structure. A local event $E$ is observed. $E$ is
consistent with a global hypothesis $H$ at a higher abstraction level. The inference from
$E$ to $H$ is made as if $E$ were specific evidence for $H$, when in fact $E$ is consistent
with many incompatible global hypotheses and provides only weak discriminating power among
them.

\begin{definition}[Abstraction level error]
An abstraction level error occurs when a claim $C$ at abstraction level $\ell_1$ is
inferred from evidence $E$ at abstraction level $\ell_2 < \ell_1$ without accounting for
the multiplicity of level-$\ell_1$ states consistent with $E$.
\end{definition}

The correct analytical response to the code leak is not a global claim about AI safety but
a local claim about pipeline constraint failure and its structural predictability. The
incident is evidence about the dynamics of high-velocity development environments under
partial constraint, not about the alignment properties of the systems being developed.
Similarly, the correct analytical response to the availability of essay-shaped generative
text is not a global claim about the death of argument but a structural claim about the
decoupling of surface from substance and the consequent shift in the reliability of
length-based signals.

The discourse on safety, in particular, is prone to this error because safety is an
inherently global property---it concerns the behavior of a system across all possible
inputs and contexts---while the available evidence is always local, consisting of specific
incidents, evaluations, and observed behaviors. The gap between local evidence and global
safety claims is not closed by accumulating more local evidence; it requires formal arguments
about the relationship between local and global properties, arguments of the kind developed
in the formal sections of this essay.

% ==========================================================
\section{Constraint Closure as Necessary but Not Sufficient}
% ==========================================================

The analysis thus far establishes constraint closure as a necessary condition for
the reliability of an output as evidence of its generating process. It does not,
however, establish sufficiency.

A system may achieve constraint closure within an incorrectly specified constraint
space. In such cases, the resulting trajectory is fully determined by the governing
structure, but that structure does not correspond to the intended domain.

\begin{definition}[Mis-specified constraint space]
A constraint space is mis-specified if it excludes trajectories that correspond
to valid solutions of the underlying problem or includes trajectories that do not.
\end{definition}

In argumentative terms, this corresponds to reasoning that is internally consistent
but based on false premises or incomplete framing. In system terms, it corresponds
to a pipeline that enforces all constraints correctly but implements the wrong
specification.

\begin{proposition}[Closure without correctness]
Constraint closure guarantees identifiability of the generating process but does
not guarantee correctness of the process with respect to an external objective.
\end{proposition}

This distinction is essential for interpreting both technical failures and
epistemological artifacts. The absence of constraint closure implies unreliability,
but its presence implies only that the process is recoverable, not that it is
correct.

The design problem is therefore twofold. First, to achieve constraint closure such
that the system's outputs are identifiable. Second, to ensure that the constraint
space itself is appropriately specified relative to the domain of interest.

Failure at the first level produces degeneracy. Failure at the second produces
coherent but incorrect systems. The two are distinct and require different forms
of analysis and remediation.

% ==========================================================
\section{Operationalization and Evaluation of Constraint Structure}
% ==========================================================

The preceding analysis establishes constraint density, identifiability, and closure
as formal properties of generative processes. To be practically useful, however,
these quantities must admit operational approximations: procedures by which a reader
or evaluator can estimate whether an observed text lies in the high-density,
identifiable regime or in the degenerate regime.

Because the generating process is not directly observable, all such procedures must
operate on the projection $Y$. The problem is therefore an inverse one: to infer
properties of the constraint structure from the observable artifact.

\subsection{Perturbation-Based Evaluation}

Let $Y$ be an argumentative text. Consider a perturbation operator $\Delta$ that
modifies a component of $Y$, such as removing a premise, altering a definition,
or reordering a segment. Let $\Delta(Y)$ denote the perturbed text.

\begin{definition}[Perturbation sensitivity]
The perturbation sensitivity of $Y$ with respect to $\Delta$ is
\[
  S_\Delta(Y) = d\bigl( \mathcal{T}(Y), \mathcal{T}(\Delta(Y)) \bigr),
\]
where $d(\cdot,\cdot)$ is a distance between trajectory sets.
\end{definition}

Intuitively, $S_\Delta(Y)$ measures how much the space of admissible trajectories
changes under a local modification of the text.

In high constraint density regimes, small perturbations induce large changes in the
admissible trajectory space, because each component of the text participates in a
network of dependencies. In degenerate regimes, perturbations have minimal effect,
because the trajectory is not tightly constrained.

\begin{proposition}[Sensitivity criterion]
If $D(Y)$ is high, then $S_\Delta(Y)$ is large for a broad class of perturbations.
If $D(Y)$ is low, then $S_\Delta(Y) \approx 0$ for most perturbations.
\end{proposition}

This provides a practical test: texts that remain coherent under arbitrary local
modification are likely to be low-density, while texts that degrade under targeted
perturbation exhibit genuine constraint structure.

\subsection{Counterfactual Trajectory Analysis}

A second operational approach is to examine counterfactual continuations. Let $Y$
be a partial text, and let $\mathcal{C}(Y)$ denote the set of admissible continuations.

\begin{definition}[Continuation entropy]
The continuation entropy of $Y$ is
\[
  H_{\text{cont}}(Y) = H(\mathcal{C}(Y)),
\]
the entropy of the distribution over admissible continuations.
\end{definition}

Low continuation entropy indicates that the trajectory is strongly constrained:
only a small number of continuations are admissible. High continuation entropy
indicates degeneracy: many continuations are equally plausible.

\begin{proposition}[Continuation constraint relation]
Constraint density is inversely related to continuation entropy. In particular,
\[
  D(Y) \uparrow \quad \Rightarrow \quad H_{\text{cont}}(Y) \downarrow.
\]
\end{proposition}

Empirically, this can be approximated by attempting to generate or enumerate
continuations and measuring their divergence. A text for which many divergent
continuations remain plausible is structurally unconstrained.

\subsection{Dependency Depth and Long-Range Consistency}

Constraint structure manifests most clearly in long-range dependencies. Let
$Y = (y_1, \dots, y_n)$ be a sequence of statements. Define a dependency relation
$y_i \rightarrow y_j$ if the admissibility of $y_j$ depends on $y_i$.

\begin{definition}[dependency depth]
The dependency depth of $Y$ is the length of the longest chain
\[
  y_{i_1} \rightarrow y_{i_2} \rightarrow \cdots \rightarrow y_{i_k}.
\]
\end{definition}

High dependency depth indicates that early constraints propagate through the
entire trajectory, enforcing global structure. Low dependency depth indicates
that statements are locally but not globally constrained.

\begin{proposition}[Depth as a proxy for constraint density]
Texts with high dependency depth exhibit higher constraint density than texts
with shallow or fragmented dependency structures.
\end{proposition}

This provides another operational signal: genuine arguments maintain consistency
over long spans, while degenerate texts exhibit only local coherence.

\subsection{Reconstruction Stability}

A final operational measure is reconstruction stability. Let $\hat{G}(Y)$ denote
an inferred generating process from the observed text.

\begin{definition}[Reconstruction stability]
A text $Y$ has high reconstruction stability if small perturbations in $Y$
produce small changes in $\hat{G}(Y)$.
\end{definition}

In high-density regimes, the constraint structure tightly determines the
generating process, so reconstruction is stable. In degenerate regimes,
small changes in the text may correspond to entirely different generating
processes, reflecting the diffuse posterior over $\mathcal{G}$.

\subsection{Interpretation}

These operational criteria do not provide exact measurements of constraint density
or identifiability. They provide observable proxies that distinguish between regimes.

A text that is sensitive to perturbation, exhibits low continuation entropy,
maintains long-range dependencies, and yields stable reconstructions is likely to
have high constraint density and to function as reliable evidence of its generating
process.

A text that is robust under arbitrary modification, admits many divergent
continuations, exhibits only local coherence, and yields unstable reconstructions
is likely to be degenerate.

The task of evaluation in high-volume generative environments is therefore not to
assess surface quality, but to estimate which regime a given artifact occupies.
Constraint structure, though not directly observable, leaves detectable signatures
in the behavior of the text under these operations.

% ==========================================================
\section{Assembly Index and Constraint-Based Complexity}
% ==========================================================

A related attempt to quantify structural complexity is provided by assembly theory,
which assigns to an object an \emph{assembly index} defined as the minimal number
of steps required to construct the object from primitive components.

Formally, for an object $O$ composed of $K$ subunits, the assembly index $a_O$
satisfies
\[
  \log_2(K) \leq a_O \leq K - 1,
\]
with lower values corresponding to structures that can be constructed through
reusable subcomponents and higher values corresponding to structures requiring
sequential construction.

Assembly theory interprets this quantity as a measure of the amount of selection
required to produce an object, and proposes that objects with high assembly index
are unlikely to arise in the absence of structured generative processes.

\subsection{Relation to Constraint Density}

The assembly index can be interpreted within the present framework as a lower bound
on the constraint structure required to generate an object. A high assembly index
implies that the generative process must traverse a constrained sequence of states,
reducing the admissible trajectory space.

In this sense, assembly index is correlated with constraint density: objects that
require many non-redundant construction steps impose stronger constraints on the
generating trajectory.

\subsection{Limitation: Absence of Projection Analysis}

However, assembly theory operates entirely at the level of generation and does not
address the inverse problem of reconstruction from observation.

In particular, the assembly index $a_O$ is a property of the minimal construction
path, not of the degeneracy of the observation-consistent set under projection.
Two objects with identical assembly indices may differ substantially in the extent
to which their generating processes are identifiable from their observable
representations.

The present framework introduces effective degeneracy $H(\config_o^{\prior})$ as
the relevant quantity for this inverse problem. While assembly index constrains
the forward generative complexity of an object, effective degeneracy determines
the recoverability of that object from its projection.

\subsection{Forward vs Inverse Complexity}

This distinction can be formalized as follows:

\begin{itemize}[leftmargin=2em]
\item Assembly index $a_O$ measures \emph{forward complexity}: the minimal effort
required to construct an object.
\item Effective degeneracy $H(\config_o^{\prior})$ measures \emph{inverse complexity}:
the uncertainty remaining about the generating process given an observation.
\end{itemize}

High forward complexity does not imply low inverse complexity. An object may be
difficult to construct yet easy to reconstruct if the prior strongly constrains
the space of admissible generators.

\subsection{Interpretation}

Assembly theory provides a useful characterization of generative difficulty, but
does not by itself determine whether an observed artifact serves as reliable
evidence of its origin. That determination requires analysis of projection,
prior structure, and degeneracy, as developed in the preceding sections.

The distinction mirrors that between effort and inference: the cost of producing
an object is not, in general, equal to the information required to identify the
process that produced it.

% ==========================================================
\section{Implications for System Design and Argumentative Practice}
% ==========================================================

The analysis yields parallel design implications for software systems and written discourse,
unified by the principle that the reliability of observable outputs as evidence of underlying
structure depends on the degree to which the generating process has achieved constraint
closure.

For AI system pipelines, the primary implication is that security cannot be achieved by
increasing the vigilance of human operators in high-velocity environments, because the
corollary of Section 7 establishes that constraint enforcement probability decreases
monotonically with velocity under human-in-the-loop verification. The appropriate
remediation is structural: replacing probabilistic human verification with architectural
invariants that are enforced by the system itself. Admissibility conditions must be embedded
in the pipeline structure rather than delegated to operator judgment. Automation, in this
context, functions as a mechanism for achieving constraint closure rather than merely for
increasing throughput.

The near-identifiability result of Section 3 has an additional implication for system design.
Because reconstruction cost is a decreasing function of prior expressiveness in the
reconstruction domain, enclosure strategies that rely on the opacity of the configuration
space become less reliable as the shared prior of the field matures. This suggests that the
long-term competitive advantage of proprietary systems lies not in the opacity of their
implementations but in properties that are genuinely difficult to reconstruct: accumulated
data, evaluation infrastructure, deployment scale, and organizational coordination. These
are high-dimensional properties whose projections remain genuinely degenerate even under
mature priors.

For argumentative practice, the implication is that the specification of constraints is
constitutive of the essay rather than preparatory to it. The premises, definitions, tensions,
and exclusions that govern an argument are not scaffolding to be removed before the essay
begins; they are the constraint structure that gives the essay its necessity. Making this
structure explicit in the text increases process legibility and thereby increases the trust
that a reader can rationally extend to the argument's conclusions.

The role of generative systems in this context is appropriately understood as a tool for
exploring constrained spaces rather than as a source of unconstrained text. When the
constraint structure of an argument is specified in advance---through explicit premises,
a defined problem, or a set of tensions to be resolved---generative systems can function
as efficient navigators of the constrained trajectory space, producing text whose necessity
is determined by the specified constraints. Absent such specification, the system produces
text that is formally correct but substantively indeterminate, and the burden of constraint
remains entirely with the reader, who cannot reconstruct a constraint structure that was
never imposed.

% ==========================================================
\section{Conclusion}
% ==========================================================

The two phenomena examined in this essay---the accidental release of proprietary AI system
code and the degeneration of argumentative essay form under generative pressure---are unified
by a common formal structure. Both arise when high-dimensional systems operate under partial
constraint, producing outputs whose surface properties substantially exceed their underlying
necessity. Both are best understood not as isolated failures but as structurally predictable
consequences of systems that have not achieved constraint closure.

The analysis supports a general principle governing the reliability of observable artifacts
as evidence of underlying generative structure: in any domain where complex processes
produce observable outputs, the reliability of those outputs as evidence depends on the
degree to which the generative process was constrained. As the cost of producing
surface-plausible outputs decreases---whether through the automation of software deployment
or the automation of text generation---the burden of signaling constraint shifts from the
artifact to the process. Artifacts that make their constraint structure legible survive as
reliable evidence. Artifacts that present surface plausibility without constraint structure
become uninformative, regardless of their surface quality.

For software systems, this principle implies that security must be architectural rather than
operational, that constraint closure must be structurally enforced rather than behaviorally
maintained, and that the effective security of enclosure is bounded above by the
reconstruction cost in the prior landscape of the reconstruction domain.

For written discourse, it implies that the essay survives as a form not by resisting the
availability of unconstrained generation but by making its constraint structure explicit
enough to remain identifiable under projection. The essay that achieves this is one in
which every transition is determined by the governing structure, every claim reduces the
space of admissible continuations, and the resulting trajectory is sufficiently constrained
that a reader can reconstruct the argumentative process from the observable text. This is
not a new standard for essays; it is the original one, newly visible because the prior
landscape of text production has changed enough to make its absence obvious.

Modern AI systems are best understood as dynamical fields subject to projection, constraint,
and diffusion rather than as static artifacts amenable to simple notions of security or
ownership. Modern argumentative texts are best understood as trajectories through constrained
spaces rather than as documents containing information. In both cases, the surface is a
projection of an underlying process, and the value of the surface as evidence depends
entirely on what can be inferred about that process from the projection.

\newpage 
\section*{Appendices}

% ==========================================================
\appendix
\section{Constraint Density and Informational Identifiability}
% ==========================================================

This appendix formalizes the notion of constraint density introduced informally in
Sections 10 through 12 and establishes its relationship to informational gain and
the identifiability of generating processes under projection.

\subsection{Constraint Density}

Let $Y$ denote an observable text and let $\mathcal{T}(Y)$ denote the set of
argumentative trajectories consistent with $Y$. Let $\mathcal{T}_0$ denote the
space of all admissible trajectories in the absence of constraints.

\begin{definition}[Constraint density]
The constraint density of $Y$ is defined as
\[
  D(Y) = 1 - \frac{H(\mathcal{T}(Y))}{H(\mathcal{T}_0)},
\]
where $H(\cdot)$ denotes the entropy of the corresponding trajectory distribution
under a reference prior over $\mathcal{T}_0$.
\end{definition}

Intuitively, $D(Y)$ measures the fraction of the original trajectory space eliminated
by the constraints encoded in $Y$. A text with $D(Y) \approx 0$ imposes few constraints
and leaves most trajectories admissible. A text with $D(Y) \approx 1$ sharply restricts
the space of admissible trajectories.

\begin{remark}
Constraint density is invariant under reparameterizations of $\mathcal{T}_0$ that
preserve entropy, and thus depends only on the relative reduction of the trajectory
space rather than its absolute representation.
\end{remark}

\subsection{Constraint Density and Mutual Information}

Let $G$ denote the generating process, modeled as a random variable taking values
in a space $\mathcal{G}$ of admissible generators. The observable text $Y$ is a
projection of $G$ through the generation map.

\begin{theorem}[Constraint density bounds mutual information]
Let $I(Y; G)$ denote the mutual information between the observable text and the
generating process. Then, under mild regularity conditions,
\[
  I(Y; G) \geq H(G) \cdot D(Y),
\]
up to a constant factor determined by the alignment between the trajectory space
$\mathcal{T}$ and the generator space $\mathcal{G}$.
\end{theorem}

\begin{proof}
The observable text $Y$ restricts the space of admissible trajectories from
$\mathcal{T}_0$ to $\mathcal{T}(Y)$. This induces a restriction on the space of
generators consistent with $Y$. The entropy reduction in the trajectory space
provides a lower bound on the entropy reduction in the generator space, since
each generator induces a trajectory. The result follows from the data processing
inequality applied to the mapping $G \mapsto \mathcal{T} \mapsto Y$.
\end{proof}

This result formalizes the intuition that constraint density is a proxy for
informational content: texts that eliminate more admissible trajectories provide
more information about the generating process.

\subsection{Identifiability Threshold}

\begin{definition}[Identifiability threshold]
Let $\epsilon > 0$. A text $Y$ is $\epsilon$-identifiable if the posterior
distribution over generators satisfies
\[
  H(G \mid Y) \leq \epsilon.
\]
\end{definition}

\begin{proposition}[Threshold condition]
There exists a threshold $D^* \in (0,1)$ such that if $D(Y) > D^*$, then $Y$
is $\epsilon$-identifiable for sufficiently small $\epsilon$.
\end{proposition}

\begin{proof}
As $D(Y) \to 1$, the entropy of $\mathcal{T}(Y)$ approaches zero, implying that
the trajectory is nearly unique. Since generators that produce distinct trajectories
must differ in their internal structure, the posterior over $G$ concentrates.
Continuity of entropy implies the existence of a threshold $D^*$ beyond which
$H(G \mid Y)$ falls below $\epsilon$.
\end{proof}

This establishes a formal criterion for when an essay functions as reliable
evidence of its generating process: it must exceed a critical constraint density.

\subsection{Degenerate Regime}

\begin{definition}[Degenerate text]
A text $Y$ is degenerate if $D(Y) \approx 0$.
\end{definition}

In this regime, the observation-consistent set over trajectories is nearly identical
to the unconstrained space. Consequently, the posterior over generators remains
diffuse, and the text provides negligible information about the process that
produced it.

\begin{corollary}
In the degenerate regime, $I(Y; G) \approx 0$, and the text is informationally
indistinguishable from a sample drawn directly from the prior over generators.
\end{corollary}

This formalizes the notion of essay-shaped text that lacks argumentative content:
it is a projection that fails to reduce uncertainty about its origin.

\subsection{Constraint Density and Length}

Finally, we relate constraint density to length.

\begin{proposition}[Length is neither necessary nor sufficient]
There exist texts $Y_1, Y_2$ such that $\text{len}(Y_1) > \text{len}(Y_2)$ but
$D(Y_1) < D(Y_2)$.
\end{proposition}

\begin{proof}
Construct $Y_1$ as a long sequence of locally plausible but unconstrained statements,
and $Y_2$ as a short sequence of tightly constrained logical steps. The entropy
reduction induced by $Y_2$ exceeds that of $Y_1$, despite its shorter length.
\end{proof}

This result formalizes the collapse of length as a reliability signal discussed in
Section 12. Constraint density, not length, determines the informational value of
a text.

\subsection{Interpretation}

The formal results of this appendix support the central thesis of the essay. Essays
that achieve high constraint density function as identifiable projections of their
generating processes and therefore serve as reliable evidence of reasoning. Essays
that fail to impose constraints remain in the degenerate regime, where their surface
structure provides no information about their origin.

In the presence of generative systems capable of producing low-density text at scale,
the distribution of observed texts shifts toward the degenerate regime. The task of
both writer and reader becomes one of restoring or detecting constraint density
sufficient to cross the identifiability threshold.

% ==========================================================
\section{Constraint Density, Kolmogorov Complexity, and MDL}
% ==========================================================

This appendix connects the notion of constraint density introduced in Appendix A
to Kolmogorov complexity and the Minimum Description Length (MDL) principle. The
goal is to relate the identifiability of argumentative trajectories to the
compressibility of the generating process.

\subsection{Kolmogorov Complexity of Generating Processes}

Let $G$ denote a generating process and $Y$ the observable text it produces. Let
$K(G)$ denote the Kolmogorov complexity of $G$, defined as the length of the shortest
program that produces $G$ on a fixed universal Turing machine.

Let $K(Y)$ denote the Kolmogorov complexity of the text. In general, $K(Y) \leq K(G) + c$,
since $Y$ can be generated by simulating $G$ and recording its output.

However, the inverse problem is of primary interest: given $Y$, what is the minimal
description length of a generator consistent with $Y$?

\begin{definition}[Minimal generating description]
The minimal generating description length of $Y$ is
\[
  K^*(Y) = \min_{G : G \mapsto Y} K(G).
\]
\end{definition}

This quantity measures the complexity of the simplest process capable of producing
the observed text.

\subsection{Constraint Density and Compression}

Constraint density can be reinterpreted as a measure of how strongly $Y$ restricts
the class of generators consistent with it. When $D(Y)$ is high, the set
\[
  \mathcal{G}_Y = \{ G : G \mapsto Y \}
\]
is small and concentrated, implying that $K^*(Y)$ is close to the true complexity
of the generating process.

When $D(Y)$ is low, $\mathcal{G}_Y$ is large and diffuse, and there exist many
simple generators that can produce $Y$. In this case, $K^*(Y)$ may be much smaller
than the complexity of the original generating process.

\begin{proposition}[Constraint density and minimal description length]
Let $Y$ be an observable text. Then $D(Y)$ is monotonically increasing in the
expected value of $K^*(Y)$ over the prior on generating processes.
\end{proposition}

\begin{proof}
As constraint density increases, the admissible set $\mathcal{G}_Y$ shrinks,
eliminating low-complexity generators that would otherwise be consistent with $Y$.
Thus the minimum of $K(G)$ over $\mathcal{G}_Y$ increases. Monotonicity follows
from the fact that adding constraints can only reduce the admissible set.
\end{proof}

This result provides a compression-theoretic interpretation of constraint density:
texts with high constraint density require more complex generators to produce them,
while low-density texts can be generated by simple, generic processes.

\subsection{Minimum Description Length Principle}

The MDL principle states that the best explanation for data $Y$ is the one that
minimizes the total description length
\[
  L(G, Y) = K(G) + K(Y \mid G),
\]
where $K(Y \mid G)$ is the conditional description length of $Y$ given $G$.

In the present context, $K(Y \mid G)$ is small when $G$ deterministically generates $Y$,
so the MDL objective reduces approximately to minimizing $K(G)$ over generators
consistent with $Y$.

\begin{definition}[MDL-optimal generator]
An MDL-optimal generator for $Y$ is
\[
  G^* = \arg\min_{G : G \mapsto Y} K(G).
\]
\end{definition}

The MDL principle therefore selects the simplest generator consistent with the
observed text. This has different consequences depending on the constraint density
of $Y$.

\subsection{Degenerate Regime Under MDL}

When $D(Y)$ is low, there exist simple generators that produce $Y$ without encoding
any meaningful constraint structure. The MDL-optimal generator in this case is a
low-complexity process that reproduces the surface form of the text without
instantiating an underlying argument.

\begin{corollary}[MDL collapse in the degenerate regime]
If $D(Y) \approx 0$, then the MDL-optimal generator $G^*$ is drawn from the class
of generic prior-based generators, and does not recover the original generating
process.
\end{corollary}

This formalizes the intuition that essay-shaped text produced without constraints
can be explained by a simple generative model, and therefore carries little
information about any deeper reasoning process.

\subsection{High-Density Regime and Generator Recovery}

When $D(Y)$ is high, the admissible generator set $\mathcal{G}_Y$ excludes simple
generators. Any generator capable of producing $Y$ must encode the constraint
structure that determines the trajectory.

\begin{theorem}[Recovery under high constraint density]
If $D(Y) > D^*$ for some threshold $D^*$, then the MDL-optimal generator $G^*$
approximates the true generating process up to equivalence under the constraint
structure.
\end{theorem}

\begin{proof}
For sufficiently high constraint density, all generators consistent with $Y$ must
encode the same constraint structure. Differences between generators are confined
to representations that are equivalent under this structure. Since MDL selects the
simplest such generator, it recovers a representative of the equivalence class of
true generators.
\end{proof}

This establishes a compression-theoretic analogue of the identifiability threshold
from Appendix A.

\subsection{Interpretation}

The connection to Kolmogorov complexity and MDL reinforces the central thesis of
the essay. The value of an argumentative text lies in its ability to constrain the
space of generators sufficiently that the simplest explanation for the text must
encode the structure of the argument itself.

In the absence of such constraints, the simplest explanation is a generic
generator that produces essay-shaped text without underlying necessity. The text
is then compressible in a way that discards any interpretation of it as a
structured argument.

Thus, constraint density serves as a bridge between information-theoretic,
computational, and epistemological perspectives. It determines not only whether a
text is informative about its generating process, but whether that process can be
recovered through compression-based inference.

% ==========================================================
\section{Assembly Index of Constraint Structures}
% ==========================================================

This appendix provides a structured mapping between the formal objects introduced
in the main text and their operational, geometric, and statistical interpretations.
Its purpose is to unify the theoretical machinery with observable diagnostics,
thereby reducing the gap between abstract analysis and empirical evaluation.

\subsection{Configuration Space Decomposition}

Let the configuration space $\config$ be decomposed as a product
\[
  \config = \config_{\text{code}} \times \config_{\text{infra}} \times \config_{\text{workflow}} \times \config_{\text{latent}}.
\]

Each component corresponds to a distinct class of constraints:

\begin{itemize}[leftmargin=2em]
\item $\config_{\text{code}}$: syntactic and semantic program structure.
\item $\config_{\text{infra}}$: deployment environment, access controls, and system topology.
\item $\config_{\text{workflow}}$: sequencing constraints and transition rules.
\item $\config_{\text{latent}}$: informal practices, conventions, and tacit knowledge.
\end{itemize}

The latent component is not directly observable but induces a prior $\prior$ over
$\config$, shaping effective degeneracy.

\subsection{Hessian Structure of Reconstruction Energy}

Recall the reconstruction energy
\[
  \energy(X) = \| \proj(X) - o \|^2 + \lambda (-\log \prior(X)).
\]

The Hessian of $\energy$ is
\[
  \nabla^2 \energy(X) = \nabla^2 \| \proj(X) - o \|^2 + \lambda \nabla^2 (-\log \prior(X)).
\]

The first term encodes projection curvature, while the second encodes prior curvature.

\begin{definition}[Condition number]
Let $\kappa$ denote the ratio of the largest to smallest eigenvalues of $\nabla^2 \energy$.
\end{definition}

\begin{proposition}[Convergence rate]
The convergence rate of reconstruction is inversely proportional to $\kappa$.
Low $\kappa$ implies rapid convergence to a prior-dominant attractor.
\end{proposition}

This formalizes the empirical observation that reconstruction is fast when the
prior strongly constrains admissible configurations.

\subsection{Diffusion on Network Topologies}

Let the information network be represented as a graph $G = (V, E)$ with Laplacian $L$.
The diffusion equation becomes
\[
  \frac{\partial \rho}{\partial t} = -D L \rho.
\]

The spectrum of $L$ determines the rate and pattern of diffusion.

\begin{proposition}[Spectral diffusion rate]
The second-smallest eigenvalue $\lambda_2$ of $L$ (the algebraic connectivity)
controls the speed of global information spread.
\end{proposition}

Highly connected networks (large $\lambda_2$) accelerate entropy diffusion,
reducing the feasibility of containment after leakage.

\subsection{Stochastic Modeling of Workflow Constraints}

Let the pipeline be modeled as a Markov chain with states $s \in S$ and transition matrix $P$.

Manual steps introduce stochastic transitions:
\[
  P(s_{t+1} \mid s_t) = p_i \cdot \delta_{\text{adm}} + (1 - p_i) \cdot \delta_{\text{escape}},
\]
where $p_i$ depends on cognitive load and iteration velocity $v$.

\begin{proposition}[Velocity-dependent failure]
If $p_i = p_i(v)$ is decreasing in $v$, then the probability of remaining within
$S_{\text{adm}}$ decays exponentially in the number of manual steps.
\end{proposition}

This provides a stochastic foundation for trajectory escape under partial closure.

\subsection{Metric for Argumentative Necessity}

Let $Y$ be a text and let $\Delta_i$ denote a perturbation of component $i$.

Define the necessity metric
\[
  N(Y) = \mathbb{E}_i \big[ d(Y, \Delta_i(Y)) \big],
\]
where $d$ measures structural degradation.

\begin{proposition}[Necessity-density relation]
$N(Y)$ is monotonically increasing in constraint density $D(Y)$.
\end{proposition}

High necessity implies that perturbations disrupt coherence, indicating strong
constraint structure. Low necessity implies robustness under modification,
indicating degeneracy.

\subsection{Unified Mapping}

The following correspondences summarize the structural equivalences developed in the essay:

\begin{center}
\begin{tabular}{ll}
\textbf{Formal Object} & \textbf{Operational Interpretation} \\
\hline
Projection $\proj$ & Observable artifact (code or text) \\
Prior $\prior$ & Shared conventions / training distribution \\
Fiber $\config_o$ & Space of admissible reconstructions \\
Entropy $H$ & Effective degeneracy \\
Energy $\energy$ & Reconstruction objective \\
Hessian $\nabla^2 \energy$ & Reconstruction stability / convergence \\
Laplacian $L$ & Network diffusion structure \\
Markov transitions $P$ & Workflow reliability \\
Constraint density $D$ & Informational clarity \\
Necessity $N$ & Fragility under perturbation \\
\end{tabular}
\end{center}

\subsection{Interpretation}

The assembly index makes explicit that the same mathematical structures govern both
technical systems and argumentative texts. Each object admits both a formal definition
and an observable proxy, allowing inference to move between levels.

The central implication is that failures of security, clarity, and trust are not
domain-specific phenomena. They arise from the same underlying condition: insufficient
constraint relative to the dimensionality of the system.

The role of the appendix is therefore not merely to collect technical details, but to
demonstrate that the theoretical framework is closed under interpretation. Every formal
object corresponds to a measurable or observable property, and every observable failure
can be traced back to a breakdown in constraint structure.

\newpage
% ==========================================================
\begin{thebibliography}{99}
% ==========================================================

\bibitem{shannon1948}
C.~E. Shannon,
``A Mathematical Theory of Communication,''
\textit{Bell System Technical Journal}, 27(3):379--423, 1948.

\bibitem{jaynes1957}
E.~T. Jaynes,
``Information Theory and Statistical Mechanics,''
\textit{Physical Review}, 106(4):620--630, 1957.

\bibitem{cover2006}
T.~M. Cover and J.~A. Thomas,
\textit{Elements of Information Theory}, 2nd ed.
Wiley-Interscience, 2006.

\bibitem{kolmogorov1965}
A.~N. Kolmogorov,
``Three Approaches to the Quantitative Definition of Information,''
\textit{Problems of Information Transmission}, 1(1):1--7, 1965.

\bibitem{grunwald2007}
P.~Gr\"unwald,
\textit{The Minimum Description Length Principle}.
MIT Press, 2007.

\bibitem{mackay2003}
D.~J.~C. MacKay,
\textit{Information Theory, Inference, and Learning Algorithms}.
Cambridge University Press, 2003.

\bibitem{tishby1999}
N.~Tishby, F.~C. Pereira, and W.~Bialek,
``The Information Bottleneck Method,''
Allerton Conference, 1999.

\bibitem{amari2000}
S.~Amari and H.~Nagaoka,
\textit{Methods of Information Geometry}.
AMS, 2000.

\bibitem{wainwright2008}
M.~J. Wainwright and M.~I. Jordan,
``Graphical Models, Exponential Families, and Variational Inference,''
\textit{FnT ML}, 2008.

\bibitem{bishop2006}
C.~M. Bishop,
\textit{Pattern Recognition and Machine Learning}.
Springer, 2006.

\bibitem{pearl1988}
J.~Pearl,
\textit{Probabilistic Reasoning in Intelligent Systems}.
Morgan Kaufmann, 1988.

\bibitem{shalev2014}
S.~Shalev-Shwartz and S.~Ben-David,
\textit{Understanding Machine Learning}.
Cambridge University Press, 2014.

\bibitem{raginsky2014}
M.~Raginsky and I.~Sason,
``Concentration of Measure Inequalities in Information Theory,''
\textit{FnT Communications}, 2014.

\bibitem{rockafellar1970}
R.~T. Rockafellar,
\textit{Convex Analysis}.
Princeton University Press, 1970.

\bibitem{boyd2004}
S.~Boyd and L.~Vandenberghe,
\textit{Convex Optimization}.
Cambridge University Press, 2004.

\bibitem{bertsekas1999}
D.~P. Bertsekas,
\textit{Nonlinear Programming}, 2nd ed.
Athena Scientific, 1999.

\bibitem{villani2009}
C.~Villani,
\textit{Optimal Transport: Old and New}.
Springer, 2009.

\bibitem{evans2010}
L.~C. Evans,
\textit{Partial Differential Equations}, 2nd ed.
AMS, 2010.

\bibitem{risken1989}
H.~Risken,
\textit{The Fokker-Planck Equation}.
Springer, 1989.

\bibitem{neal1993}
R.~M. Neal,
``Probabilistic Inference Using MCMC,''
Technical Report, 1993.

\bibitem{hastings1970}
W.~K. Hastings,
``Monte Carlo Sampling Methods Using Markov Chains,''
\textit{Biometrika}, 1970.

\bibitem{asmussen2007}
S.~Asmussen and P.~W. Glynn,
\textit{Stochastic Simulation}.
Springer, 2007.

\bibitem{lovasz1993}
L.~Lov\'asz,
``Random Walks on Graphs: A Survey,''
1993.

\bibitem{newman2010}
M.~E.~J. Newman,
\textit{Networks: An Introduction}.
Oxford University Press, 2010.

\bibitem{shalizi2013}
C.~R. Shalizi,
``Advanced Data Analysis from an Elementary Point of View,''
2013 (draft).

\bibitem{cronin2021}
L.~Cronin, S.~I. Walker, et al.,
``Assembly Theory Explains and Quantifies Selection and Evolution,''
\textit{Nature}, 2021.

\bibitem{cronin2017}
L.~Cronin,
``Assembly Theory,''
University of Glasgow, 2017.

\bibitem{benner2023}
S.~A. Benner,
``On the Limitations of Assembly Theory,''
2023.

\bibitem{marshall2024}
C.~Marshall et al.,
``Assembly Theory and Algorithmic Information: A Critical Analysis,''
\textit{Journal of Molecular Evolution}, 2024.

\bibitem{royalsoc2024}
A.~Group et al.,
``Abiotic Generation of High Assembly Index Structures,''
\textit{J. Royal Society Interface}, 2024.

\end{thebibliography}

\end{document}
